\documentclass[10pt]{article}
\usepackage[margin=0.9in,headheight=14pt]{geometry}
\usepackage{amsmath,amssymb}
\usepackage[T1]{fontenc}
\usepackage[utf8]{inputenc}
\usepackage{lmodern}
\usepackage{microtype}
\usepackage{booktabs,array}
\usepackage{enumitem}
\usepackage{xcolor}
\usepackage{fancyhdr}
\usepackage[hidelinks]{hyperref}
\hypersetup{pdftitle={RED TEAM EXERCISE: the swoop paper, written to be taken apart}}
\definecolor{rt}{RGB}{170,20,20}
\pagestyle{fancy}\fancyhf{}\renewcommand{\headrulewidth}{0.4pt}
\fancyhead[C]{\footnotesize\bfseries\color{rt}RED-TEAM EXERCISE. FICTIONAL AUTHOR. NOT A DERIVATION. NOT A FRAMEWORK RESULT.}
\fancyfoot[C]{\footnotesize\thepage}
\setlist{nosep,leftmargin=1.3em}
\setlength{\parskip}{4pt}\setlength{\parindent}{0pt}
\newcommand{\tell}[1]{\textsuperscript{\color{rt}\bfseries[T#1]}}
\sloppy
\begin{document}

\fbox{\parbox{0.97\linewidth}{\small\textbf{\color{rt}Read this first.} This is a deliberate exercise. Part~I is written in the voice of a fictional, ethically careless defender of the standard model who wants to take a dark energy relation for $a_0$ and call it old news. It does not derive that relation. Its numbers come from AI-run toy scripts that no outside group has reproduced, chosen to flatter it. Red marks \tell{n} point to Part~II, which shows what each passage hides. No real person is alleged to be doing this, and no evidence of anyone doing so was found. Do not circulate Part~I without Part~II.}}

\vspace{8pt}
\part*{Part I: the paper as the priest would write it}

\begin{center}
{\Large\bfseries On the Natural Emergence of a Characteristic Acceleration\\ in the Concordance Cosmology}\\[6pt]
{\small Prof.\ A.~N.~Orthodox (fictional), Institute for Concordance Studies (fictional)}
\end{center}

\begin{abstract}\noindent
The persistent numerical proximity of the galactic acceleration scale $a_0$ to cosmological quantities has long invited speculation. We demonstrate\tell{1} that no such speculation is required. In the concordance cold dark matter cosmology, the central acceleration of the $L_\star$ halo is $9.17\times10^{-11}$\,m\,s$^{-2}$, within 2\,\% of $\tfrac12 c\sqrt{G\rho_\Lambda}=9.36\times10^{-11}$\,m\,s$^{-2}$.\tell{2} A semi-empirical population of discs in such haloes reproduces a tight relation of the observed form with a characteristic scale of $9.1$--$9.2\times10^{-11}$\,m\,s$^{-2}$.\tell{3} The scale responds positively to the vacuum density through the growth history,\tell{4} supplying the long-sought origin of the apparent connection. The same haloes remove the mass deficit of galaxy clusters\tell{5} and accommodate a rising apparent scale at high redshift.\tell{6} We conclude that relations of this type are a natural corollary of the standard model and carry no independent content.\tell{7}
\end{abstract}

\section{Introduction}
It has been known since Milgrom \cite{M83} that $a_0$ lies near $cH_0$, and Milgrom \cite{M99} and Limbach, Psaltis \& \"Ozel \cite{LPO} have associated it with the vacuum. Such associations recur periodically, of late in the unrefereed grey literature, where one occasionally encounters them dressed with a specific coefficient.\tell{8} Against this, the emergence of a tight radial acceleration relation in hierarchical galaxy formation is by now a mature result \cite{N17,L17,KW17,D19}. As Ludlow et al.\ showed, it is a natural outcome of simulations of galaxy formation in cold dark matter \cite{L17}. It has remained only to show that the \emph{value} follows too.

\section{The characteristic acceleration of the $L_\star$ halo}
For an NFW halo the central acceleration $g_0=2\pi G\rho_s r_s$ is independent of the galaxy it hosts and thus constitutes the most physically fundamental\tell{9} characteristic scale. We adopt the standard concentration--mass relation (Diemer \& Joyce 2019) at $\Omega_\Lambda=0.685$, $H_0=67.4$. At the characteristic halo mass $M_{200}=2\times10^{12}\,M_\odot$ we find
\[
g_0=9.17\times10^{-11}\ \text{m\,s}^{-2}.
\]
After exploring a range of reasonable definitions,\tell{10} we regard this as the robust result. The mild mass dependence, $d\ln g_0/d\ln M\simeq0.18$, reflects the near-universal character of cold dark matter haloes.\tell{11}

\section{A semi-empirical disc population}
We populate haloes with exponential discs through abundance matching \cite{Mo13}, with realistic scatter of 0.11, 0.15 and 0.2\,dex in concentration, stellar mass and size. Taking disc sizes at the upper end of plausible size--halo scalings, and noting that feedback is widely understood to offset baryonic contraction,\tell{12} we fit the observed functional form at 0.5--5 disc scale lengths. Independent realisations give $a=9.09$, $9.19$ and $9.11\times10^{-11}$\,m\,s$^{-2}$, with a stack rms of 0.073--0.078\,dex, in pleasing accord with observation.\tell{13}

\section{The vacuum connection explained}
Expressed in units of $c\sqrt{G\rho_\Lambda}$ at the concordance parameters, the fitted scale gives $\hat\kappa=0.500$.\tell{14} The relation is therefore no mystery. The vacuum density shapes the growth history, and with it halo concentrations, so that $d\ln a/d\ln\rho_\Lambda>0$ at fixed matter density and primordial amplitude.\tell{4} A correlation of the kind sometimes attributed to new physics is thus precisely what the concordance model predicts.

\section{Clusters, the Solar System, and high redshift}
The concordance model encounters none of the embarrassments of modified dynamics. The Solar System is untouched. At $R_{500}$ the halo mass exceeds that of the modified relation by factors of 1.49--1.89 from $2\times10^{14}$ to $10^{15}\,M_\odot$, in the direction the data demand.\tell{5} An ultra-faint dwarf in a dark halo shows a dispersion above the modified relation, as observed.\tell{15}

MUSE-DARK III \cite{C26} reports a rise to $a_0(z\sim1)=2.38\times10^{-10}$\,m\,s$^{-2}$, confirmed by Alcock's analysis \cite{Al26},\tell{16} with a constant scale excluded at high significance. Hydrodynamical work finds a similar rise \cite{Ma23}. Our contracted models give 1.26, 1.50 and 2.03 at $z=1$, 1.5 and 2.33.\tell{6} Any law pinned to a constant vacuum density is thereby excluded.

\section{Conclusions}
The value of $a_0$, its proximity to the vacuum scale and its evolution all follow from standard physics with no additional parameters.\tell{17} We see no need to pursue alternatives further.

\part*{Part II: taking it apart}
\section*{What each mark hides}
\small
\begin{tabular}{@{}l>{\raggedright\arraybackslash}p{0.91\linewidth}@{}}
\toprule
\textbf{Mark} & \textbf{What the scripts actually show} \\ \midrule
T1 & ``Demonstrate'' is false. Nothing here is a derivation. \\
T2 & 2\,\% is the number after choosing the radius rule, mass, concentration relation and weighting. $g_0$ runs from $3.43\times10^{-11}$ to $2.86\times10^{-10}$ over $10^{10}$--$10^{15}\,M_\odot$. Weighted over discs it is $7.45\times10^{-11}$ by stellar mass and $5.21\times10^{-11}$ by number. Both miss. \\
T3 & Lands only with no contraction and disc size 1.345 times the base scaling. Contraction gives 2.58 times the target (Blumenthal) or 2.06 (Gnedin). Without scatter, 0.81 times. Size changes of $\pm20$\,\% give 0.91 and 1.09. About 6--7 choices. \\
T4 & \textbf{The central dodge.} A positive sign is not the relation. The slope with $\rho_{\rm DE}$ is $+0.13$ for $g_0$ and $+0.07$ for the tuned stack, against $+0.50$ for the law. Multiplying $\rho_{\rm DE}$ by 4 should double $a_0$; these barely move it. \\
T5 & The direction is right but the size isn't: 1.49--1.89 against the 2.93 measured for X-COP. Key-by-key grading was not run. The cost goes unstated: the apparent scale stops being universal. \\
T6 & Two different models are mixed. The tuned stack that gives the headline number gives 0.99, 1.07, 1.28. Only the contracted model rises, and it misses the headline value by more than a factor of 2. No single model does both. \\
T7 & ``Carry no independent content'' doesn't follow. A similar number at one mass isn't the same statement as a mass-independent constant tied to $\rho_{\rm DE}$. \\
T8 & \textbf{The credit move.} Cite only the shared ancestors, wave at the ``grey literature'', and give no DOI. This is the one step that would actually cost you credit, and it takes one sentence. Often it would be plain ignorance, not misconduct. \\
T9 & $g_0$ is the $r\to0$ limit, where baryons dominate. It is not a transition scale. ``Most fundamental'' is chosen because it lands. \\
T10 & 714 definitions were tried; 20 (2.8\,\%) land within 10\,\%, and half of those are $g_0$. For arbitrary targets across $10^{-12}$ to $10^{-9}$ the median is 12 hits, and 37\,\% of targets do as well or better. The grid is correlated, so this is indicative only. \\
T11 & A slope of 0.18 over 2 dex is not near-universal. The fitted scale runs $6.4\times10^{-11}$, $7.0\times10^{-11}$, $8.5\times10^{-11}$, $1.6\times10^{-10}$ from $10^{10}$ to $5\times10^{12}\,M_\odot$. \\
T12 & A continuous knob and a removed process, described as a modest assumption. \\
T13 & No comparison was made. SPARC's rms in the law is 0.0541\,dex. Per-galaxy $a$ scatters by 0.24\,dex and trends with surface brightness ($-0.34$ per dex at fixed speed), where the law is flat. Metrics differ, so this is indicative. \\
T14 & $\hat\kappa$ is not a constant. It is 0.432 at $\Omega_\Lambda=0.70$, 0.726 at $H_0=73$, and 0.994 at $(0.65, 73)$. The 0.500 was tuned in, and only that row is quoted. These swings are large because $\omega_m$ and $\sigma_8$ move at fixed $A_s$. \\
T15 & Ratio 1.59--1.73 between script runs, for an assumed $10^9\,M_\odot$ halo. That is accommodation, not a fit. IC~1613 goes unmentioned because it gets no rescue: rival over law is 0.94--0.97, and the law already misses it low. \\
T16 & ``Confirmed'' is false. Alcock reanalyses the same four bins, stated statistical errors only, and the paper gives the mass caveat itself. The rise travels with the mass route, which the priest leaves out. \\
T17 & Abundance matching, $c(M)$, sizes and feedback are all inputs. ``No additional parameters'' is simply untrue. Lensing also goes unmentioned: at $10^{12}\,M_\odot$ the one-halo term falls $-0.22$\,dex below the law at 300\,kpc and $-0.73$ at 1\,Mpc. The two-halo term is omitted and would need grading. \\
\bottomrule
\end{tabular}
\normalsize

\section*{Route by route}
\small
\begin{tabular}{@{}llll>{\raggedright\arraybackslash}p{0.27\linewidth}@{}}
\toprule
Route & Scale at $10^{12}\,M_\odot$ & Slope with $\rho_{\rm DE}$ & Choices & Verdict \\ \midrule
Top-hat collapse with $\Lambda$ & $4.7\times10^{-13}$--$3.1\times10^{-12}$ & 0 to $+0.46$ & -- & Can't land; needs $10^{16}$--$10^{19}\,M_\odot$ \\
NFW $g_0$ & $8.2\times10^{-11}$ & $+0.13$ & 4--5 & Lands near $2\times10^{12}$ only \\
Contraction + matching & about 2.5 times target & $+0.19$ & misses & Misses \\
Same, tuned & $9.36\times10^{-11}$ & $+0.07$ & 6--7 & Lands; slope wrong \\
Coupled dark sector & $9.36\times10^{-11}$ & $+0.50$ & 2--3 & By hand; $\alpha^3=0.0997$; not plain $\Lambda$CDM \\
\bottomrule
\end{tabular}
\normalsize

The coupled sector is the superfluid form $a_0=\alpha^3\Lambda^2/M_{\rm Pl}$ \cite{BFK}, with $\Lambda^4=\rho_{\rm DE}$ imposed. It has the law's form because the form was imposed, and $\alpha$ is far from that paper's $\alpha^3$ of about 185.

\section*{All the knobs: how perfect it can be made}
Targets and ranges were frozen before any number (\texttt{FROZEN\_TUNING.md}). The base model is Blumenthal contraction on 8 halo masses, with eleven knobs: $c(M)$ norm and slope, four SHMR parameters, contraction fraction, disc size norm and slope, gas norm and slope. Tier~1 ranges are my judgement, not a cited compilation.

\textbf{The first run found a loophole.} With ranges tripled, it reached rms 0.0000\,dex, spread 0.001\,dex and $a$ within 0.05\,\% of the law. It got there by shrinking a $10^{10}\,M_\odot$ galaxy's disc to 0.04\,kpc, so every point sat at $g_{\rm bar}$ of 93--246 times $a_0$ even at its outermost radius. There, the inner NFW acceleration is nearly constant and mimics $g_{\rm bar}+a_0$, with $g_{\rm DM}$ at 0.98\,$a_0$ for every mass. The match never touched the low-acceleration regime. Amendment A1, made after seeing this and disclosed as such, extends each fit out to $g_{\rm bar}=0.03\,a_0$.

\small
\begin{tabular}{@{}l>{\raggedright\arraybackslash}p{0.25\linewidth}cccccc@{}}
\toprule
Tier (A1 fit range) & Added freedom & T1 level & T2 flat & T3 shape & T4 slope & T5 $\kappa$ & T6 $z$ \\ \midrule
Base & none & F & F & F & F & F & F \\
1 & eleven knobs, plausible & P & F & F & F & F & F \\
2 & ranges tripled & P & P & P & F & F & F \\
3 & plus $c\propto\rho_{\rm DE}^{0.30}$ by hand & P & P & P & P & F & F \\
3b & plus a $c$ shift per cosmology and per $z$ & P & P & P & P & solved & solved \\
4 & free-form halo copied from the law & built in & built in & built in & built in & built in & built in \\
\bottomrule
\end{tabular}
\normalsize

\begin{itemize}
\item \textbf{Tier 1} reaches ratio 1.002 but has 0.247\,dex mass spread and rms 0.0226\,dex, with most knobs at their bounds. In plausible ranges, the law's level is reachable and its mass independence and shape are not, on this toy.
\item \textbf{Tier 2} gets rms 0.0054\,dex and spread 0.012\,dex, the near-perfect fit you'd expect to be asked for. Its price is discs of 0.03\,kpc at $10^{10}\,M_\odot$ and 0.14\,kpc at about $1.5\times10^{11}$, far too small to be real galaxies, plus $\log M_1$ shifted by 0.54 and $\beta$ by $-0.83$. Its slope stays at $+0.182$.
\item \textbf{Tier 3} gets the $+0.500$ slope only by writing $\rho_{\rm DE}^{0.30}$ into the concentration. Nothing computes that exponent. $\kappa$ still runs 0.455, 0.665, 0.800, and $a(z)/a(0)$ is 1.29, 1.55, 2.15, outside the band.
\item \textbf{Tier 3b} needs $c$ shifts of $+0.041$, $-0.098$, $-0.183$\,dex by cosmology and $-0.158$, $-0.247$, $-0.500$\,dex at $z$ of 1, 1.5, 2.33. Each was solved alone and the shifts were not run jointly, so ``solved'' means one at a time. The shift at $z=2.33$ is a factor of about 3 in concentration, with no process behind it. (Corrected 2026-11-02: an earlier version of this page gave 1.79, 2.07, 2.68 for $a(z)/a(0)$, $-0.306$, $-0.384$, $-0.617$\,dex and a factor of 4, from a code error that computed $c(M,z)$ in the wrong cosmology. The verdicts did not change.)
\item \textbf{Tier 4} is exact because it's the law renamed: $M_{\rm DM}=r^2(g_{\rm law}-g_{\rm bar})/G$, the phantom halo. Density stays positive at the masses checked, and its mass against NFW at $R_{200}$ is 1.56, 1.44, 1.25, 0.55. It tends to $\rho\propto r^{-2}$ outward with no process to produce that. It would also inherit every failure of the law: IC~1613, X-COP and the rest. So perfect agreement with the law is not agreement with the data.
\end{itemize}

\textbf{Reading.} A perfect match is available, and the cost is visible: impossible discs, then $\rho_{\rm DE}$ inserted by hand, then a separate knob per cosmology and redshift, then the law itself as the halo. This shows how much freedom there is, not that either side is right. It also gives any reader a checklist: disc sizes, fit range reaching $g_{\rm bar}\ll a_0$, slope with $\rho_{\rm DE}$, $\kappa$ across cosmologies.

\section*{What the exercise shows}
\begin{itemize}
\item \textbf{The rival can reach the number, not the law.} No route produced the $+\tfrac12$ slope from dynamics. So the question to ask of any such paper is its slope $d\ln a/d\ln\rho_{\rm DE}$, its $\hat\kappa$ across cosmologies, and how many definitions it tried.
\item \textbf{That is no win for the law.} $\kappa=\tfrac12$ is itself fitted, and SPARC doesn't single it out: the free fit corresponds to $\kappa$ 0.575 and $cH_0/2\pi$ to 0.557. The rival's failure doesn't make the tie real.
\item \textbf{The rival is really better in places.} It has clusters and dwarfs moving the right way and no Solar System problem. The failures are complementary, and neither absorbs the other. This says nothing about which the data favour.
\item \textbf{The honest attack needs none of this.} Cite Milgrom and Limbach, plot the law as one more $a_0(z)$ curve, and point to MUSE-DARK. That is ordinary and legitimate, and the law's real exposure.
\item \textbf{What the priest does hurts only through omission.} T8 is the credit risk, and arXiv posting, a refereed paper and direct contact address it.
\end{itemize}

\section*{Limits}
Toy models only, no simulations. Concentration and growth come from analytic relations through \texttt{colossus}; Moster 2013's parameters, the Gnedin coefficients and the splashback factor were entered from memory and aren't verified. Parameter slopes use fixed $\omega_m$, $\omega_b$, $A_s$ unless named. Rerunning the scripts reproduced the reported values except the contracted stack (2.46--2.58 times target across runs) and the dwarf ratio above. All of this was run by AI subagents and reproduced by no outside group. Literature details come from abstracts and summaries, so check them before relying on them. $\kappa=\tfrac12$ is fitted. The theory is not closed. Nothing here says the data favour this framework or $\Lambda$CDM. Scripts and outputs: \texttt{scripts/}. Number check: \texttt{audit.py}.

\begin{thebibliography}{99}\small
\bibitem{M83} M.~Milgrom, ApJ 270, 365 (1983).
\bibitem{M99} M.~Milgrom, Phys.\ Lett.\ A 253, 273 (1999).
\bibitem{LPO} C.~Limbach, D.~Psaltis, F.~\"Ozel, arXiv:0809.2790 (2008).
\bibitem{N17} J.~F.~Navarro et al., arXiv:1612.06329 (2017).
\bibitem{L17} A.~D.~Ludlow et al., PRL 118, 161103 (2017), arXiv:1610.07663.
\bibitem{KW17} B.~W.~Keller, J.~W.~Wadsley, ApJL 835, L17 (2017), arXiv:1610.06183.
\bibitem{D19} A.~A.~Dutton et al., arXiv:1902.06751 (2019).
\bibitem{Mo13} B.~P.~Moster et al., arXiv:1205.5807 (2013).
\bibitem{C26} B.~I.~Ciocan et al., A\&A 709, L16 (2026), arXiv:2604.22613.
\bibitem{Al26} G.~T.~Alcock, MNRAS 551, L1 (2026), arXiv:2609.38796.
\bibitem{Ma23} L.~Mayer et al., arXiv:2206.04333.
\bibitem{BFK} L.~Berezhiani, B.~Famaey, J.~Khoury, arXiv:1711.05748 (2018).
\end{thebibliography}
\end{document}
